% 09-use.tex: держатели, доверяющие стороны, брокер и путь постороннего.
\section{Использование удостоверения: держатели, доверяющие стороны и вход, не оставляющий записи}
\label{sec:use}

\motif{Быть проверенным не значит быть выслеженным.}

\analogy{Горожане носят свои бумаги сами. Купец узнаёт то, что купцу нужно, и не более, из бумаг, которые покупатель решил показать, а два купца, сверив свои книги, не могут понять, что обслужили одного и того же человека.}

\subsection{Кошелёк есть протокол, а не продукт}

Эталонный агент держателя, \code{сценарии/\allowbreak{}полярис-кошелёк.пи}, держит удостоверение в виде файла,
показывает его, проверяет автономно через отделяемый верификатор, предъявляет его, доказывает
принадлежность с нулевым разглашением против опубликованной эпохи, подписывает документ от имени
держателя, делегирует ограниченные полномочия агенту (раздел~\ref{sec:agents}) и входит к
доверяющей стороне. Всякое сообщение, которое он выпускает, описано так, что выпустить такое же
способен любой клиент. Родные мобильные и настольные клиенты и браузерный мост суть продуктовая
инженерия, за которую эталонная реализация на условных данных не берётся; они следовали бы
протоколу и никогда бы его не вели. Путь физического предъявления, карта, подписывающая вызов
устройства, есть эмулятор и устройство проверки из раздела~\ref{sec:core}. \sImplMark{} для протокола
и эталонного сценария; \sFut{} для родных клиентов.

Предъявление есть неподписанная обёртка вокруг подписанных вещей: подписанное эмитентом
удостоверение, необязательно подшитое утверждение о состоянии, чтобы полномочия были разрешимы без
всякой связи, необязательное доказательство принадлежности, необязательная подпись держателя над
одноразовым числом верификатора, контекст и уровень раскрытия, под которыми предъявляют, и
непрозрачный код предъявления. Верификатор решает автономно и сообщает код как присутствующий или
отсутствующий, ни разу его не истолковывая, потому что код принуждения сопоставляет лишь эмитент,
вне поля зрения; предъявление под принуждением побайтово той же формы, что и обычное. Для передачи
по куар или НФС предъявление сжимается и разбивается на кадры, каждый из которых называет дайджест
целого; получатель берёт их в любом порядке и отвергает, ещё ничего не разобрав, кадры с разными
дайджестами, пропущенный номер, противоречащий дубликат или собранную нагрузку, дайджест которой не
сошёлся. Декодер ограничен в ресурсах: самое большее 9999 кадров, 512~КиБ в сжатом виде и 2~МиБ в
развёрнутом, так что бомба распаковки отвергается на пределе, а не распаковывается. Границы были
добавлены после того, как внешнее ревью назвало этот пробел. \sImpl

\subsection{Доверяющая сторона}

Эталонная доверяющая сторона (\code{сценарии/\allowbreak{}полярис-доверяющая-сторона.пи}) решает, принять или
отвергнуть предъявление, сочетая автономную подлинность с сетевым состоянием, отвергает отозванное,
подменённое или чужое по эмитенту удостоверение и остаётся слепа к принуждению, а сквозное учение
прогоняет всю матрицу от кошелька до доверяющей стороны под настоящими подписями при каждом
выпуске. Доверяющая сторона регистрируется у эмитента как организация с областью действия, которую
ограничение схемы сужает ровно до \code{проверка}, \code{опознать} или обоих, так что четвёртый
род доступа есть конституционное изменение, а не настроечное. Личность не становится входным
продуктом по недосмотру. Доверяющая сторона есть также предмет, о котором
записываются решения: регистрация, отключение ступени с нулевым разглашением, которую обязаны были
проходить её держатели, отключение или повторное включение, каждое пишет событие в режиме
добавления, изменение, ослабляющее то, чему сторона обязана удовлетворять, помечается как таковое
и отвергается без изложенной причины, а изменение, сделанное из консоли базы данных, записывается
на тех же условиях, что и сделанное через инструмент, потому что пишет триггер, а не инструмент.
Первая версия правила, решающего, что считать ослаблением, ранжировала возможность по длине её
имени, так что выдача возможности, которой сторона не имела, читалась как безобидная; это нашли,
прогнав правило против загруженной базы данных, а не прочитав его. \sImpl

\subsection{Вход без записи о входе}

Брокер входа позволяет доверяющей стороне предложить \emph{вход по удостоверению Полярис} через
код авторизации с ПКСЕ. Держатель предъявляет удостоверение в экземпляре своего выдавшего органа по
владению, а доверяющая сторона обменивает полученный код на идентификационный токен, подписанный
выдавшим ведомством, который проверяет автономно. Четыре стража не дают этому стать записью о
входе. Субъект токена есть подескриптор доверяющей стороны из раздела~\ref{sec:privacy}, никогда не
идентификатор человека и не одно и то же значение у всякой доверяющей
стороны. Токен не несёт утверждений сверх словаря: контекст, уровень раскрытия, достигнутый уровень
достоверности, состояние регистрации, моменты времени. Маршрут авторизации не пишет ничего, а
единственная запись конечной точки токена есть хеш кода в регистр, который не держит ни субъекта,
ни доверяющей стороны. И принуждение обслуживается неотличимо. \sImpl

Требования доверяющей стороны, обязательная ступень с нулевым разглашением или обязательное
состояние регистрации, связывают орган и не передаются лишь в запросе со стороны держателя:
политика регистрируется в строке доверяющей стороны и применяется первой, запрос вправе добавить
требование и никогда не снять, а контекст, отличный от зарегистрированного, отвергается. Код
авторизации не только подписан, но и зашифрован. А у контракта соответствия позже обнаружилось, что он проверяет
подпись идентификационного токена и никогда его аудиторию или одноразовое число, так что
верификатор, не обращавший внимания ни на то ни на другое и потому принимавший токен, вычеканенный
для одной доверяющей стороны, у другой, проходил каждый случай; четыре случая это исправили.
\sImpl

\subsection{Путь постороннего}

Читатель способен дойти до одного принятого предъявления программами, которых здесь никто не писал, на чистой машине,
примерно за десять минут и не клонируя репозитория: поставить \code{полярис-оид4вп} из реестра,
запустить опубликованный образ кошелька по закреплённому дайджесту, выработать ключи верификатора,
выдать кошельку удостоверение, связанное с ключом, который он выработал сам, запустить верификатор,
предъявить. Страницу, объясняющую как, исполнили из пустого каталога прежде, чем записали, и по
дороге она поймала собственный дефект (буферизованный канал, оставлявший журнал верификатора
пустым), и заканчивается она той строкой, которая важна:
\code{<- 200 подлинно, утверждения ['пк', 'фамилия', 'имя', 'выдан', 'эмитент', 'тип\_уд']}. Если у
читателя она не сработает, говорит страница, это будет полезнейшая ошибка, какую проект способен
получить. Никто за пределами репозитория о прохождении этого пути пока не сообщал. \sImplMark{} для
пути; \sExt{} для проходящего, который не есть автор.

\subsection{Что намеренно не построено}

У органа нет сессии, нет экрана согласия, нет документа обнаружения и нет хореографии браузерных
перенаправлений: доверяющая сторона интегрируется по спецификации формата передачи и по реестру.
Браузерный мост ВебАутн не построен; ключ на стороне держателя, на который он опирался бы, теперь
существует. Ни одна доверяющая сторона за пределами репозитория не интегрировалась с родным
протоколом. Один сторонний кошелёк предъявил по ОИД4ВП, и ни одна сторонняя доверяющая сторона
ещё не запускала принявшего его верификатора. \sExt

\begin{layercard}{Карточка слоя: использование}
\qa{Что это}{Протокол кошелька и эталонный сценарий, предъявления и разбиение на куар-кадры, эталонная доверяющая сторона, регистрация доверяющей стороны с ограниченной областью и записываемой историей, брокер входа и путь постороннего до принятого предъявления из стороннего кошелька.}
\qa{Зачем оно существует}{Удостоверение, которое некому держать, предъявлять и принимать, есть строка в базе данных. Этот слой есть место, где удостоверение встречает программы, которых эмитент не запускает.}
\qa{Чему доверяет}{Подписанным объектам внутри предъявления; ключу эмитента ради идентификационного токена; зарегистрированной политике доверяющей стороны.}
\qa{Чему отказывает}{Истолкованию кода предъявления; записи о входе; области вне ограничения ПРОВЕРКА; попытке запроса держателя ослабить зарегистрированное требование доверяющей стороны; понижению планки доверяющей стороны без причины в записи; распаковке неограниченной нагрузки.}
\qa{Что видит}{Доверяющая сторона видит значение удостоверения и вердикты, а на пути с нулевым разглашением доказательство и дескриптор с областью. Орган видит доказательство владения и поглощает хеш кода.}
\qa{Чего не видит}{Орган не видит, кто куда вошёл; доверяющая сторона не видит записи человека.}
\qa{Если откажет}{Переигранный код натыкается на регистр поглощённых; неверное предъявление, невыполненная ступень, предъявитель только с правом проверки у брокера, устаревший или подменённый кадр отвергаются каждый в своём учении; ослабленная планка доверяющей стороны без причины отвергается базой данных.}
\qa{Сейчас или потом}{Сейчас, как протокол, и сейчас для одного стороннего кошелька на одном пути. Родные клиенты и сторонняя доверяющая сторона потом.}
\qa{Внутри и снаружи}{Внутри: граница приватности. Снаружи: федерация, потому что доверяющая сторона редко доверяет одному-единственному органу.}
\end{layercard}

\nextlayer{Доверяющая сторона, доверяющая одному органу, есть дело простое. У общества органов
много, и ни один из них не должен уметь говорить за остальных, и ни один не должен держать запись
о каждом человеке. Следующий слой есть то, как независимые органы признают друг друга, не
схлопываясь в один.}
